\scanpage{31}
$\widehat A'$ en un schéma formel abélien $\widehat A$.

Le th. de comparaison alg. formel nous ramène à prouver que $A$ est algébrique.

On est ainsi ramené à prouver le

\uline{Th.} Soit $R$ anneau loc noeth complet, $S=\operatorname{Spec}R$, $\widehat A$ un schéma \inserted{abélien formel} abélien sur $\widehat S=\operatorname{Spf}\unsure{B|R}$. Conditions équivalentes

(i) $\widehat A$ \unsure{algébrique}, i.e. $\exists$ schéma abélien $A$ sur $S$, tel que $\widehat A$ soit isom. à son complété formel le long de $A_s$.

(ii) $\exists$ morphisme propre \uline{surjectif} \deleted{\unsure{surjectif}} $f:S'\longrightarrow S$, \deleted{\unsure{et}} tel que l'image inverse $\widehat A'$ de $\widehat A$ sur $\widehat S'$ soit algébrique.

(iii) Pour toute \deleted{algèbre finie $B$ sur $A$} quotient \deleted{fini} \inserted{intègre} $B$ de $A$, de dim $1$, si $B'$ est le normalisé de $B$, et \deleted{$\unsure{\widehat B}$} $\widehat T'=\unsure{\operatorname{Specf}|\operatorname{Spf}}B'$, --- $\widehat A_{\widehat T'}$ est algébrique.

\margin{\uline{NB} La démonstration \unsure{montre} qu'on peut se borner aux $x\in S$, $\dim\{\unsure{\bar x|x}\}=1$, qui ne sont pas contenus dans les $\bigcup_i T_i$, où les $T_i$ sont des \unsure{fermés} dans $S$ de \unsure{toute} famille \uline{dénombrable}.}

\deleted{N.B. On peut \unsure{démontrer} le th. si $\dim R=1$, et en utilisant \unsure{l'algébricité} de l'espace des \unsure{faisceaux}.}

\uline{Dém.} $(i)\Longrightarrow(ii)$ est trivial.

$(ii)\Longrightarrow(iii)$ Car introduisant $T'=S'\times_S T$, on voit que $\widehat T'\mathbin{-}\widehat T$ ``rend $\widehat A_{\widehat T}$ algébrique''. Or on peut trouver, comme $T$ est un trait complet, un trait $T_1$ \deleted{et plat} \inserted{(fini sur $T$)}, et un $T$-morphisme $T_1\mathbin{-}T'$. Donc $\widehat A_{\widehat T_1}$ est algébrique, donc définit un schéma abélien $A_{T_1}$ sur $T_1$. Ce dernier est muni d'une donnée de descente rel à $T_1\mathbin{-}T$ [grâce au th. de comparaison], il suffit \deleted{\unsure{encore}} de voir que la donnée est effective, ce qui est évident, car $T_1\mathbin{-}T$ est \inserted{fid} plat fini, et $A_1$ est projectif.

\margin{\deleted{Pour \unsure{utiliser} \unsure{cela}, la \unsure{N.Q.|N.O.} \unsure{sur} \illegible{1} \unsure{implique} \illegible{1} \unsure{limite} \unsure{imposer} une famille \unsure{indénombrable} \unsure{à éviter}.}}

\mspage{30}
